% CAV submission — Springer LNCS format.
% Requires llncs.cls (Springer's LNCS class) and splncs04.bst.
% Download: https://www.springer.com/gp/computer-science/lncs/conference-proceedings-guidelines
% Compile: pdflatex main; bibtex main; pdflatex main; pdflatex main
%
% CAV is double-blind at submission: keep authors/institute anonymized until camera-ready.
\documentclass[runningheads]{llncs}

\usepackage[T1]{fontenc}
\usepackage{amsmath,amssymb,amsfonts,mathtools}
\usepackage{booktabs}
\usepackage{xcolor}
\usepackage{listings}
\usepackage{tikz}
\usetikzlibrary{arrows.meta,positioning}
\usepackage[hidelinks]{hyperref}

% ---- notation ----
\newcommand{\R}{\mathbb{R}}
\newcommand{\Q}{\mathbb{Q}}
\newcommand{\Z}{\mathbb{Z}}
\newcommand{\net}{\mathcal{N}}
\newcommand{\reg}{R}
\newcommand{\dom}{\mathcal{B}}          % input box
\newcommand{\leaf}{L}
\newcommand{\leaves}{\mathcal{L}}
\newcommand{\vipr}{\textsc{Vipr}}
\newcommand{\scip}{\textsc{SCIP}}
\newcommand{\tool}{\textsc{AptpCheck}}
\newcommand{\pytool}{\textsc{APTPchecker}}
\newcommand{\lean}{\textsc{Lean\,4}}
\newcommand{\mathlib}{\textsc{Mathlib}}
\newcommand{\sound}{\textsf{certified}}
\newcommand{\relu}{\mathrm{ReLU}}
\newcommand{\code}[1]{\texttt{#1}}

\lstdefinestyle{lean}{
  basicstyle=\ttfamily\footnotesize,
  keywordstyle=\bfseries,
  commentstyle=\itshape\color{gray!70!black},
  showstringspaces=false,
  columns=fullflexible,
  breaklines=true,
  keywords={theorem,lemma,def,structure,inductive,by,fun,forall,exists,if,then,else,match,with,let,intro,exact,simp,have,Prop,Bool},
}
\lstset{style=lean}

\begin{document}

\title{Kernel-Checked Verification of Neural-Network\\ Activation-Pattern-Tree Proofs with Exact\\ Solver Certificates}
\titlerunning{Kernel-Checked Verification of Neural-Network Proofs}

\author{Anonymous Author(s)}                 % camera-ready: real names
\institute{Anonymous Institution}            % camera-ready: real affiliation
\authorrunning{Anonymous}

\maketitle

\begin{abstract}
Neural-network verification (NNV) tools have matured to the point that they
prove nontrivial properties of realistic networks, but the tools themselves are
large, highly optimized, and floating-point based, and thus may be unsound. The
recently proposed \emph{Activation-Pattern-Tree Proof} (APTP) format lets a NNV
tool emit a verifier-independent proof that an external \emph{checker} can
validate, moving trust from the (complex) verifier to the (simpler) checker.
We ask the natural next question: \emph{who checks the checker?} We present
\tool, a checker for APTP proofs whose soundness argument is \emph{formalized} in
\lean{}/\mathlib{} and whose core reasoning steps are checked by the logical kernel,
with \emph{no} appeal to floating point and \emph{no} trust in any external solver.
\tool{} computes over exact rationals; it discharges each proof leaf by re-checking a
\vipr{} certificate emitted by an \emph{untrusted} run of the exact solver \scip{},
and it establishes that the proof leaves \emph{cover} the whole input region. We
factor the end-to-end guarantee into three obligations---per-leaf refutation, region
coverage, and encoding soundness---and prove that their conjunction implies the
original property holds for \emph{every} input in the box. Mechanized and
\emph{axiom-clean} so far: encoding soundness for \emph{arbitrary-depth} feed-forward
ReLU networks (the full big-$M$ layer induction), a \emph{complete, automatic,
runnable} proof-by-cases certificate checker over \vipr{} (Farkas / rounding / integer
splits), the region-coverage check, and---moving them out of the trusted base---both
input-file parsers, proved by round-trip. The exact evaluator is validated bit-for-bit
against an independent oracle. We describe the formalization, the trusted computing
base, and the integration with \scip/\vipr.
\keywords{Neural network verification \and Proof checking \and Interactive
theorem proving \and Lean \and Exact arithmetic \and Farkas certificates \and
Mixed-integer programming.}
\end{abstract}

\section{Introduction}\label{sec:intro}

Neural networks are increasingly deployed in settings where a wrong output has a
real cost, and a substantial body of work now \emph{verifies} that a network
satisfies a formal specification---e.g., that a small perturbation of any input in
a region cannot change the classification, or that a controller never commands an
unsafe action~\cite{katz2017reluplex,wang2021betacrown,vnncomp2023}. The winning
tools of the annual VNN-COMP competition combine bound propagation, abstract
interpretation, and branch-and-bound (BaB) over neuron activation states, and are
among the most sophisticated pieces of software in formal methods.

That sophistication is also a liability. These tools are tens of thousands of
lines of GPU-accelerated, floating-point code; a subtle rounding bug or an
incorrect abstract transformer can make a tool report ``verified'' for a network
that in fact violates the property. Because a verifier's answer is a single bit,
such unsoundness is invisible. The standard remedy in SAT/SMT and mixed-integer
programming is \emph{proof production}: the tool emits a certificate that an
independent, much simpler \emph{checker} validates, so trust shifts from the tool
to the checker and the certificate
format~\cite{isac2022marabou,cheung2017vipr,eifler2022exact}.

For NNV, the \pytool{} line of work introduced the
\emph{Activation-Pattern-Tree Proof} (APTP) format: a \emph{verifier-independent}
proof that a network satisfies a linear property over an input box. An APTP proof
is a set of \emph{leaves}, each a partial \relu{} activation pattern; the leaves
partition the input region, and each leaf comes with the claim that no input
consistent with that pattern violates the property. The associated checker builds
a mixed-integer linear program (MILP) per leaf and calls an off-the-shelf solver
to discharge it. This is a real improvement---the checker is far simpler than the
verifier---but two trust gaps remain. First, the checker is itself an unverified
Python program using floating-point pre/post-processing. Second, and more
fundamentally, it \emph{trusts a commercial MILP solver} (Gurobi): the solver's
``infeasible'' is taken on faith, reintroducing exactly the kind of large,
floating-point, closed-source dependency the proof format was meant to eliminate.

\paragraph{Contributions.} We close both gaps. We present \tool, a checker for
APTP proofs whose soundness argument is \emph{formalized in \lean{}~4} against
\mathlib{} and whose core reasoning steps are checked by the \lean{} kernel with no
use of \code{native\_decide}, no floating point, and no trusted solver. Concretely:

\begin{itemize}
  \item We give a precise \lean{} statement of \emph{what an accepted APTP
    proof guarantees}, and factor it into three obligations---\emph{per-leaf
    refutation}, \emph{coverage}, and \emph{encoding soundness}---with a top-level
    theorem that their conjunction implies the property holds for every input in
    the box (\S\ref{sec:checking}).
  \item We remove the trusted solver. We run the \emph{exact} solver \scip{}
    (\scip-Ex/SoPlex) as an \emph{untrusted oracle} and re-check its \vipr{}
    certificate~\cite{cheung2017vipr,eifler2022exact} inside \lean{} over exact
    rationals; SCIP leaves the trusted computing base entirely. We mechanize a
    \emph{complete, automatic, runnable} proof-by-cases checker over the three \vipr{}
    step rules (Farkas combination, integer rounding, integer split), proved sound and
    demonstrated running on a parsed certificate (\S\ref{sec:certificates}).
  \item We identify and correctly handle two subtleties that the original checker
    glosses: the propositional \emph{coverage} check is a \emph{sound
    over-approximation} that is only correct up to a global polarity symmetry, and
    each leaf is itself a \emph{MILP} (not an LP) because partial patterns leave
    unstable neurons free (\S\ref{sec:checking}--\ref{sec:certificates}).
  \item We isolate \emph{encoding soundness}---that the rational constraint system
    over-approximates the network's real behaviour on each region---as the central
    theorem, and \emph{prove it for arbitrary-depth} feed-forward \relu{} MLPs: both
    cornerstones (big-$M$ ReLU soundness, interval-bound validity) and the full
    layer-induction assembly (\S\ref{sec:formalization}).
  \item We \emph{verify both input parsers}---the exact network file and the
    proof-tree file---by round-trip against printers, moving them out of the trusted
    base; this required rewriting the parsers' \code{partial}/imperative loops as
    total structural recursions without changing their behaviour
    (\S\ref{sec:formalization}).
  \item We describe the exact-rational substrate, the bit-exact float-to-rational
    weight decoding, and the trusted computing base
    (\S\ref{sec:formalization}--\ref{sec:implementation}), and evaluate on APTP
    proofs generated by state-of-the-art verifiers (\S\ref{sec:eval}).
\end{itemize}

To our knowledge \tool{} is the first NNV proof checker to reduce its soundness to a
proof-assistant kernel: the numerically hard \emph{search} is delegated to an
untrusted MILP solver (\scip), whose \vipr{} proof certificate is then
\emph{re-checked in the kernel} over exact rationals rather than trusted.

\section{Background}\label{sec:background}

\paragraph{Networks and properties.}
We consider feed-forward networks $\net : \R^{n_0}\!\to\R^{n_\ell}$ that alternate
affine layers $x \mapsto Wx+b$ (dense or convolutional) with the elementwise
$\relu(t)=\max(t,0)$, and optionally reshape (\code{Flatten}) layers. A
\emph{verification query} is a triple $(\net, \dom, (c, \rho))$ where
$\dom=[\underline x,\overline x]\subseteq\R^{n_0}$ is an input box and
$(c,\rho)\in\R^{n_\ell}\times\R$ is a linear property. The query \emph{holds} iff
\begin{equation}\label{eq:property}
  \forall x\in\dom.\quad c\cdot \net(x) \;>\; \rho .
\end{equation}
For example, ``output $0$ always exceeds output $1$'' is $c=(1,-1,0,\dots)$,
$\rho=0$. Conjunctive/disjunctive specifications reduce to a family of such rows,
each checked independently.

\paragraph{Activation patterns.}
Each hidden \relu{} neuron $k$ has a pre-activation value; its \emph{sign} at input
$x$ is \emph{active} ($\ge 0$) or \emph{inactive} ($<0$). Fixing the sign of every
neuron makes $\net$ \emph{affine}, which is why NNV tools branch on neuron signs. A
\emph{partial activation pattern} (a \emph{leaf}) is a finite set of signed neuron
indices $\pm k$: $+k$ means neuron $k$ is active, $-k$ inactive. A leaf $\leaf$
defines a region
\begin{equation}\label{eq:region}
  \reg(\leaf) \;=\; \{\, x\in\dom : \text{the sign of every neuron mentioned in }\leaf\text{ at }x\text{ agrees with }\leaf \,\}.
\end{equation}

\paragraph{MILP encoding.}
A query restricted to a leaf is encoded as a MILP over rationals: affine layers
become equalities $v = Wu+b$; a \relu{} neuron that is provably stable on $\dom$
becomes an identity or a constant $0$; an \emph{unstable} neuron becomes a fresh
output variable with a binary indicator $a\in\{0,1\}$ and the four big-$M$
constraints
\[
  v \ge u,\quad v \ge 0,\quad v \le u - \underline u\,(1-a),\quad v \le \overline u\, a,
\]
where $[\underline u,\overline u]$ are valid pre-activation bounds. The property is
folded into the last layer ($W_{\text{last}}\!\leftarrow c\,W_{\text{last}}$), so a
single terminal variable equals $c\cdot\net(x)$; minimizing it and proving the
minimum positive discharges \eqref{eq:property} on that leaf.

\paragraph{APTP proofs.}
An APTP proof (Fig.~\ref{fig:tree}) is a set of leaves $\leaves$, serialized in an
SMT-LIB-like syntax that declares the input/output variables and the piecewise-linear
(\relu) neurons, asserts the input box and the property, and asserts a
disjunction-of-conjunctions (\emph{DNF}) over neuron signs---one conjunction per
leaf. The intended reading: the leaves cover $\dom$, and each is refuted. The
checker's job is to validate both.

\begin{figure}[t]
\centering
\begin{tikzpicture}[
  every node/.style={font=\footnotesize}, level distance=9mm, sibling distance=22mm,
  edge from parent/.style={draw,-{Latex[length=1.4mm]}}]
\node (r) {$\net,\dom,(c,\rho)$}
  child {node {$N_4^-$} }
  child {node {$N_4^{+}$}
    child {node {$N_2^-$}}
    child {node {$N_2^{+}$}
      child {node {$N_1^{+}$}}
      child {node {$N_1^-$}}
    }
  };
\end{tikzpicture}
\caption{An APTP proof as a tree of \relu{} sign splits (the running
\code{sample.aptp}). Each leaf is a partial pattern, e.g.\
$\{N_2^{+},N_1^{-},N_4^{+}\}=\{2,-1,4\}$; the proof never branches on $N_3$, which
therefore stays a free (unstable) binary in each leaf's MILP. Leaves tile $\dom$;
each carries a refutation certificate.}
\label{fig:tree}
\end{figure}

\paragraph{Exact solving and \vipr{} certificates.}
Floating-point MILP solvers can err on both feasibility and optimality. The
\emph{exact} solver \scip-Ex uses rational LP (SoPlex) with safe bounding and,
crucially, can emit a \vipr{} certificate---\emph{Verifying Integer Programming
Results}~\cite{cheung2017vipr,eifler2022exact}---recording every reasoning step.
A \vipr{} file lists variables, constraints, the \emph{relation to prove} (RTP:
infeasibility or an objective bound), primal solutions, and a sequence of
\emph{derived} constraints, each justified by one of: \code{asm} (a branching
assumption), \code{lin} (a sign-consistent nonnegative rational combination of
earlier constraints---an LP dual / Farkas step), \code{rnd} (a \code{lin} followed
by integer rounding), or \code{uns} (combining the two children of an integer split
$x_i\le k \lor x_i\ge k{+}1$, discharging the split assumption). A certificate is
accepted when the final derived constraint is an absurdity (for infeasibility) or
dominates the claimed bound, with an \emph{empty assumption set}. Checking a \vipr{}
file requires only \emph{exact rational arithmetic} and the \emph{easy} direction of
Farkas' lemma; it never re-solves.

\section{A worked example}\label{sec:example}

We instantiate \emph{every} step of the pipeline on the toy network
\code{sample.onnx} and proof \code{sample.aptp}, and for each step we name the \lean{}
object that computes it \emph{and} the lemma that discharges it, sketching the proof.
All arithmetic is exact; decimals such as $0.2$ are shown for readability but the
checker carries the exact binary32 value $\mathtt{0x3E4CCCCD}=13421773/67108864$.

\paragraph{The network.} A $2$-input, $2$-output ReLU MLP
$\mathrm{Lin}\!\to\!\relu\!\to\!\mathrm{Lin}\!\to\!\relu\!\to\!\mathrm{Lin}$ with
\[
W_0=\begin{psmallmatrix}-0.5&0.2\\ 0.3&0.2\end{psmallmatrix},\,
b_0=\begin{psmallmatrix}-0.5\\ 0.3\end{psmallmatrix};\;
W_2=\begin{psmallmatrix}-0.9&0.6\\ 0.6&0.9\end{psmallmatrix},\,
b_2=\begin{psmallmatrix}0.2\\ -0.1\end{psmallmatrix};\;
W_4=\begin{psmallmatrix}-0.1&0.2\\ -0.3&0.5\end{psmallmatrix},\,
b_4=\begin{psmallmatrix}-0.4\\ -0.5\end{psmallmatrix}.
\]
The first \relu{} layer supplies hidden neurons $N_1,N_2$, the second $N_3,N_4$
(1-based, layer-major, matching the reference variable numbering). Write
$p^{(1)}=W_0x+b_0$, $h^{(1)}=\relu(p^{(1)})$, $p^{(2)}=W_2h^{(1)}+b_2$,
$h^{(2)}=\relu(p^{(2)})$, $y=W_4h^{(2)}+b_4$; so $N_1,N_2$ are the signs of
$p^{(1)}$ and $N_3,N_4$ of $p^{(2)}$.

\paragraph{The query.} Input box $\dom:\ x_0\in[-2,2],\,x_1\in[-1,1]$; the property
\code{(assert (<= Y\_0 Y\_1))} gives the objective row $c=(1,-1)$, $\rho=0$, so the
goal is $\forall x\in\dom.\ y_0-y_1>0$ (i.e.\ $Y_0>Y_1$). Our parser returns exactly
this ($\dom$, $c$, $\rho$) and the leaves below.

\paragraph{Step 1 --- exact weight decoding (\code{float32ToRat}).}
Every weight in the \code{.net} file is stored as its raw 32-bit hex pattern (e.g.\
$0.2 \to \mathtt{0x3E4CCCCD}$). Our decoder extracts sign, exponent, and significand
directly from the bits and returns the exact fraction --- $13421773/67108864$ for
$0.2$. No decimal approximation is involved, so the checker's internal representation
is \emph{exactly} the number the original network file encodes. This is why the whole
pipeline can be exact: we never start from a rounded value.

\emph{What we proved.} The decode is exact by definition (it is a closed formula over
the bit-fields). We validate it by computing: $\mathtt{0x3F800000}\mapsto1$,
$\mathtt{0xC0000000}\mapsto-2$, etc.\ --- matching the IEEE-754 spec.

\paragraph{Step 2 --- verified parsing (\code{parseNet}, \code{parseAptp}).}
Two parsers read the network file and the proof-tree file. We proved that each parser
is ``faithful'': if you take a valid in-memory representation and print it back to
text, re-parsing that text gives back the \emph{same} object. Formally:
\code{parseNet(printNet(r)) = Ok(rawToNet(r))}, and likewise for \code{parseAptp}.

\emph{Why this matters.} It means the parsers cannot silently eat or corrupt data. If
the file is well-formed, the parser reconstructs exactly what the file represents. We
proved this by showing each sub-parser (numbers, token lists, variable declarations,
box bounds, each leaf) inverts its matching printer, then composed them.

On the sample, parsing produces $\dom=[-2,2]\times[-1,1]$, $c=(1,-1)$, $\rho=0$,
neurons $\{1,2,3,4\}$, leaves $\{[-4],[-2,4],[2,1,4],[2,-1,4]\}$.

\paragraph{Step 3 --- running the network exactly, and bridging to the proof model.}
The parsed network runs as ordinary code: multiply arrays, add biases, apply ReLU. We
also have a mathematical model (inductive \code{MLP}) on which the soundness theorems
are stated. \code{toMLP\_eval} proves the two \emph{compute the same outputs} for all
inputs.

\emph{Why.} The soundness theorems say ``the encoding matches the MLP.'' The bridge
says ``the MLP matches the parsed network.'' Chaining them: soundness applies to the
actual parsed network, not just a separate mathematical copy.

On the sample, \code{toMLP} succeeds (the network is a valid MLP with $\text{inDim}=
\text{outD}=2$) and exact evaluation gives e.g.\ $\net(1,1)_0 =
-25991905743211293818487/75557863725914323419136$.

\paragraph{Step 4 --- bounding each neuron's range (interval analysis).}
For each neuron, we compute the widest output it can produce given the input box. The
rule: ``positive weights times the low end + negative weights times the high end =
lower bound'' (and vice versa for the upper bound). On the sample:
\[
N_1\!\in\![-1.7,0.7],\quad N_2\!\in\![-0.5,1.1],\quad
N_3\!\in\![-0.43,0.86],\quad N_4\!\in\![-0.1,1.31].
\]
All four neurons can be either positive or negative (``unstable''), so none can be
simplified away.

\emph{What we proved (\code{affine\_interval\_sound}).} For any input $x$ inside the
box, the neuron's true value is between the computed bounds. The argument: each term
$w\cdot x_i$ is sandwiched between $w\cdot\mathrm{lo}_i$ and $w\cdot\mathrm{hi}_i$
(positive coefficient: multiply by the lower bound gives the lower bound; negative:
it flips); summing all the terms preserves the sandwich.

\paragraph{The proof tree.} \code{sample.aptp} asserts a four-leaf DNF; our parser
yields $\leaves=\{[-4],\,[-2,4],\,[2,1,4],\,[2,-1,4]\}$ over the signs of
$N_1,N_2,N_4$:
\[
\leaf_1:\lnot N_4,\quad \leaf_2:\lnot N_2\wedge N_4,\quad
\leaf_3:N_2\wedge N_1\wedge N_4,\quad \leaf_4:N_2\wedge\lnot N_1\wedge N_4.
\]
The proof never branches on $N_3$, which therefore stays a free binary neuron inside
every leaf's MILP---a concrete reason the leaves are MILPs, not LPs.

\paragraph{Step 6 --- coverage: the leaves cover everything.}
At any input $x$, each neuron is either ``on'' ($\ge0$) or ``off'' ($<0$). The
branched neurons here are $(N_1,N_2,N_4)$, giving $2^3=8$ possible combinations. The
four leaves together must cover all eight. We enumerate them: $N_4$ off
$\Rightarrow\leaf_1$; $N_4$ on, $N_2$ off $\Rightarrow\leaf_2$; $N_4$ on, $N_2$ on,
$N_1$ on $\Rightarrow\leaf_3$; $N_4$ on, $N_2$ on, $N_1$ off $\Rightarrow\leaf_4$.
Every combination is accounted for.

\emph{What we proved (\code{coverage\_sound}).} \code{checkCoverage} lists all $2^k$
on/off combinations of the branched neurons, then checks each one is matched by some
leaf. If the check passes, then no matter what the real neurons do, some leaf
applies. Proved by induction: the enumeration is complete (every real assignment agrees
with some enumerated one on the relevant neurons), and our matching test is correct.
$N_3$ never appears in any leaf, so it does not affect coverage---both its on/off
states are handled inside each leaf.

\paragraph{Step 7 --- encoding a leaf as a system of inequalities.}
Take leaf $\leaf_3=[2,1,4]$: neurons $N_2,N_1,N_4$ are ``on'' (their values are
$\ge0$), while $N_3$ is unbranched (could be on or off). The encoder writes
inequalities that any real network execution in this leaf \emph{must} satisfy:
The encoder is \emph{leaf-aware}: a neuron the leaf pins ``on'' needs no $0/1$ helper
at all --- its output simply \emph{is} its input. Variables are laid out in fixed
slots (inputs $v_0,v_1$; layer-1 pre-activations $v_2,v_3$, helpers $v_4,v_5$, outputs
$v_6,v_7$; layer-2 pre-activations $v_8,v_9$, helper $v_{10}$, outputs $v_{12},v_{13}$;
network outputs $v_{14},v_{15}$):
\begin{itemize}
\item \textbf{Box:} the inputs are in range ($v_0\le2$, $-v_0\le2$, $v_1\le1$, $-v_1\le1$).
\item \textbf{Affine equalities:} each pre-activation equals the linear function of its
      input, e.g.\ $v_2=-0.5v_0+0.2v_1-0.5$.
\item \textbf{Folded ReLUs for the pinned $N_1,N_2,N_4$:} just $v_6=v_2$, $v_7=v_3$,
      $v_{13}=v_9$ (``output $=$ input'') --- no helper, no integer variable.
\item \textbf{Big-$M$ ReLU for the unbranched $N_3$ only:} four inequalities using the
      one $0/1$ helper $v_{10}$ to say ``$v_{12}=\max(v_8,0)$'', with the range from
      Step~4 ($[-0.43,0.86]$) as the ``big-$M$'' constants.
\item \textbf{Sign constraints:} the leaf says $N_1,N_2,N_4\ge0$, so $v_2\ge0$,
      $v_3\ge0$, $v_9\ge0$.
\item \textbf{Negated property:} the objective $v_{14}-v_{15}\le\rho$ (the opposite of
      what we want to prove).
\end{itemize}
In total: 31 rows + the objective row over 16 variable slots, with a \emph{single}
integer variable ($v_{10}$) --- because three of the four neurons are pinned by the
leaf. (Across the four leaves the integer count is $3,2,1,1$: exactly the unfixed
neurons of each.)

\emph{What we proved (\code{encFold\_overapprox}).} If you plug in the
\emph{true} network values (the actual $x$, the true pre/post values of each neuron,
the true output), then \emph{every single row above is satisfied}, and the objective
form evaluates to $c\cdot\net(x)$. This is proved by induction over the layers: each
layer's rows hold because they are literally restating what the layer computes (affine)
or what ReLU means (the four big-$M$ rows encode ``output = max(input, 0)''). The
essential point: the encoding is a \emph{relaxation} --- the real execution is always
inside the feasible set, but there may be ``extra'' fake solutions too. This means the
encoding can never be over-tight (which would be unsound).

\paragraph{Step 8 --- proving the leaf has no counterexample (certificate replay).}
The system from Step~7 (the encoding rows plus the negated property) has no solution
--- meaning every point that satisfies the encoding must \emph{violate} the negated
property, i.e.\ $c\cdot\net(x)>\rho$. This is what we want, but we need proof.

The external solver (exact \scip{}) reports ``infeasible'' and provides a certificate
--- a list of derived constraints, each built from earlier ones by three kinds of
steps:
\begin{enumerate}
\item \textbf{Add (lin):} multiply some existing rows by nonneg.\ weights and sum them.
\item \textbf{Round (rnd):} if a row's left-hand side is always an integer, round the
      right-hand side down.
\item \textbf{Case-split (uns):} since $v_7$ is $0$ or $1$, split into two sub-problems
      and refute each.
\end{enumerate}
The final derived row is $0\le-1$ (a contradiction), showing no solution exists.

\emph{What we proved (\code{checkVipr\_sound}).} Our verified checker reads the
certificate directly (no untrusted intermediate conversion) and re-does every step in
exact arithmetic. At the end it checks that the contradiction has no remaining
assumptions. The soundness theorem: if the checker accepts, then genuinely no
assignment satisfies the rows. The proof goes by induction over the derivation list,
using three lemmas: ``add'' is valid by the linearity of $\le$; ``round'' is valid
because an integer $\le$ a fraction implies it $\le$ the floor; ``case-split'' is valid
because an integer is either $\le k$ or $\ge k{+}1$.

Concretely, for the small example $2x\ge1\wedge2x\le1$ with $x$ integer, the
certificate splits $x$ at $0$: in the $x\le0$ branch, $2\!\cdot\!(x\le0)+1\!\cdot\!
(2x\ge1)\Rightarrow 0\le-1$; in the $x\ge1$ branch, $2\!\cdot\!(x\ge1)+1\!\cdot\!
(2x\le1)\Rightarrow 0\le-1$. Both branches reach a contradiction, so the system is
infeasible --- our checker confirms (\code{checkVipr} returns \code{true}).

\paragraph{Step 9 --- putting it together.}
The logic chain:
\begin{enumerate}
\item Step~8 says the rows (from Step~7) together with $c\cdot y\le\rho$ have no solution.
\item Step~7 says the real network execution \emph{does} satisfy those rows.
\item Therefore the real execution must \emph{violate} $c\cdot y\le\rho$, i.e.\
      $c\cdot\net(x)>\rho$ on this leaf.
\end{enumerate}
Repeat for all four leaves. Step~6 says every input belongs to some leaf. Together:
$\forall x\in\dom,\ c\cdot\net(x)>\rho$ (Theorem~\ref{thm:main}). Finally, Step~3
shows this applies to the actual parsed network (not just a separate mathematical
copy), because the two compute the same thing.

This whole chain is one machine-checked theorem, \code{certify\_network\_sound}: if the
per-leaf check \code{leafCheck} passes for every leaf and the coverage check passes, then
$\forall x\in\dom,\ c\cdot\net(x)>\rho$ on the \emph{parsed} network. \code{leafCheck}
bundles three decidable conditions: (i) the verified certificate checker accepts
(\code{checkSem}); (ii) \emph{certificate--encoding correspondence}: every constraint row
in the certificate's \code{CON} section semantically matches a row of this leaf's
encoding (Step~7's rows, plus the binaries' $0\le b\le1$ bounds) --- ``semantically''
because the solver renumbers variables (undone, untrusted, from the certificate's own
variable names --- a wrong renaming can only cause rejection) and reorders terms (handled
by the proven order-insensitive form equality); and (iii) every variable the certificate
declares integer is one of the encoder's ReLU indicators. Without (ii) the certificate
would prove infeasibility of \emph{its own} rows, not necessarily \emph{ours}; with it,
the guarantee attaches to the network.

\paragraph{Verdict.} $\sound\iff$ coverage $\wedge$ every leaf refuted; composing
(Theorem~\ref{thm:main}) gives $Y_0>Y_1$ on all of $\dom$. The property is tight:
at the corner $x=(2,1)$ the exact output is $y\approx(-0.308,-0.313)$, so
$y_0-y_1\approx+0.005>0$---close enough that sampling could not establish it, but a
proof does.

\paragraph{Step 10 --- the real run (not a mock-up).} We ran the assembled tool
\code{aptpcheck} on this exact \code{sample.net}/\code{sample.aptp} pair end-to-end
against \emph{official} exact \scip{} (the \code{scipoptsuite-10.0.1} release, built with
its exact rational mode and the exact \textsc{SoPlex} LP solver). For each leaf the tool
emits the MILP as an \code{.mps} file with \emph{exact rational} coefficients (e.g.\
$0.2$ is written and read back losslessly as $13421773/67108864$; a fused coefficient
appears verbatim as \code{V0 R4 1/2}), declares \emph{only} the unfixed neurons' ReLU
indicators as integer, and invokes exact \scip{}, which returns a \vipr{} certificate.
The verified \code{leafCheck} (Step~9) then validates each certificate \emph{and} its
correspondence to the leaf's encoding. The console output:
\begin{center}\small
\begin{tabular}{l}
\texttt{Coverage: OK (leaves tile the activation-pattern cube)}\\
\texttt{Objective 0: c=[1, -1], rhs=0}\\
\texttt{~~leaf 0: CERTIFIED (verified leafCheck: cert validated + matches this leaf's encoding)}\\
\texttt{~~leaf 1: CERTIFIED \dots\quad leaf 2: CERTIFIED\quad leaf 3: CERTIFIED}\\
\texttt{CERTIFIED: $\forall$ x $\in$ dom, c $\cdot$ net(x) > $\rho$}\\
\end{tabular}
\end{center}
The four certificates \scip{} produced are real solver output, not hand-written: over
$32$--$38$ constraints and $16$ variables per leaf they run $97$, $273$, $177$ and $983$
derivation steps respectively; two leaves additionally use Gomory rounding cuts
(\code{rnd} steps) alongside the signed-multiplier combinations (\code{lin}). The
leaf-aware encoding is directly visible in the integer-variable counts: leaves
$\leaf_1,\dots,\leaf_4$ declare $3,2,1,1$ integers --- exactly the neurons each leaf
leaves unfixed. All four certificates are accepted by the kernel-checked checker
(\S\ref{sec:certificates}), so the tool's verdict is precisely the hypothesis of
\code{certify\_network\_sound}. On this network the leaf-aware encoding keeps each MILP
small enough that exact \scip{} proves every leaf with signed combinations and Gomory
cuts alone --- no branching --- which is why these certificates use only \code{lin} and
\code{rnd}. To exercise branching too, we ran exact \scip{} on the integer-infeasible
knapsack $6x+10y+15z=23$ (bounded integers; LP-feasible), which it solved by
branch-and-bound, emitting a $32$-step certificate that uses two assumptions
(\code{asm}), a case-split (\code{uns}) and eight Gomory cuts (\code{rnd}) alongside
signed combinations (\code{lin}). Our verified checker accepts it --- and correctly
\emph{rejects} it if the integrality declarations are stripped. This exercises the full
chain, branching included, on genuine untrusted solver certificates.

\section{What checking an APTP proof means}\label{sec:checking}

We first pin down the guarantee and its decomposition; \S\ref{sec:certificates}
and \S\ref{sec:formalization} discharge the pieces.

\begin{definition}[Accepted proof]
$\tool(\net,\dom,(c,\rho),\leaves,\mathit{certs})$ returns \sound{} iff (i) every
leaf $\leaf\in\leaves$ has a \vipr{} certificate in $\mathit{certs}$ that our
checker validates, and (ii) the coverage check on $\leaves$ succeeds.
\end{definition}

\begin{theorem}[Soundness]\label{thm:main}
If $\tool(\net,\dom,(c,\rho),\leaves,\mathit{certs})=\sound$ then
$\forall x\in\dom.\ c\cdot\net(x) > \rho$.
\end{theorem}

The proof factors through three lemmas, each proved separately in \lean.

\begin{lemma}[Per-leaf refutation]\label{lem:leaf}
If our checker validates the certificate for leaf $\leaf$, then
$\forall x\in \reg(\leaf).\ c\cdot\net(x) > \rho$.
\end{lemma}

\begin{lemma}[Coverage]\label{lem:cover}
If the coverage check on $\leaves$ succeeds, then
$\forall x\in\dom.\ \exists \leaf\in\leaves.\ x\in\reg(\leaf)$.
\end{lemma}

\begin{proof}[of Theorem~\ref{thm:main}]
Fix $x\in\dom$. By Lemma~\ref{lem:cover} there is $\leaf\in\leaves$ with
$x\in\reg(\leaf)$; by Lemma~\ref{lem:leaf} for that $\leaf$,
$c\cdot\net(x)>\rho$. \qed
\end{proof}

This is exactly the informal ``the checker visits all nodes of the BaB tree'':
coverage (Lemma~\ref{lem:cover}) is that the leaves \emph{tile} the region, and
per-leaf refutation (Lemma~\ref{lem:leaf}) is that every tile is discharged. The
optimized driver of \pytool{} additionally \emph{subsumes} a leaf by any solved
more-general leaf; subsumption is sound because refuting a larger region refutes
its sub-regions, and the driver terminates because its queue strictly shrinks each
round. We prove subsumption sound but do not otherwise need the queue: soundness
requires only a certificate per (unsubsumed) leaf plus coverage.

\paragraph{The coverage check, precisely.}\label{par:coverage}
The original checker builds a CNF whose clauses are the leaves' literal-sets and
asserts UNSAT. This decides ``the DNF $\bigvee_{\leaf}\bigwedge_{\ell\in\leaf}\ell$
is a tautology''---but only over \emph{free independent booleans} $N_k$, and only
because flipping the polarity of \emph{every} literal is a satisfiability-preserving
bijection, so $\mathrm{UNSAT}\!\left(\bigwedge_\leaf\bigvee_\ell \ell\right)
=\mathrm{UNSAT}\!\left(\bigwedge_\leaf\bigvee_\ell \lnot\ell\right)$, the latter being
exactly ``the negation of the DNF is unsatisfiable.'' We formalize the clean version:
negate every literal and check the CNF is UNSAT. Two points matter for soundness.
\emph{(a)} A tautology over the free boolean cube is \emph{sufficient} for coverage
(every real sign vector is \emph{some} assignment, hence satisfies some leaf) though
not necessary---the check may reject a proof that only covers \emph{reachable}
patterns; this costs completeness, never soundness. \emph{(b)} A neuron the proof
never branches on (e.g.\ $N_3$ in Fig.~\ref{fig:tree}, whose exact range
$[-0.43,0.86]$ straddles $0$) is simply left free: both its signs lie inside each
region, so coverage over the branched neurons suffices while the unbranched neuron
stays a binary inside each leaf's MILP. Separately, if a producer \emph{folds} a
neuron it claims stable, we re-verify that stability with exact interval bounds (part
of encoding soundness, \S\ref{sec:formalization}).

\paragraph{Leaves are MILPs, not LPs.}
Because leaves are \emph{partial}, a neuron unstable on $\dom$ but not mentioned in
$\leaf$ keeps its binary indicator inside $\reg(\leaf)$'s encoding. Hence a per-leaf
certificate is a full \vipr{} derivation (with \code{rnd} and \code{uns}), not a single
Farkas multiplier vector. (Requiring \emph{complete} patterns would make each leaf
affine and its certificate a single \code{lin} step; we support that as an easier mode
but do not assume it.)

\section{Removing the trusted solver}\label{sec:certificates}

Lemma~\ref{lem:leaf} is where the numerically hard work lives, and where the original
checker trusts Gurobi. We keep the hard work in a solver but trust \emph{none} of it.

\paragraph{Encoding.} For a leaf $\leaf$ we build a rational system
$A x \le b$ (variables $x$ collecting inputs and per-layer pre/post-activations) as in
\S\ref{sec:background}, add the leaf's sign constraints, and add the \emph{negated}
property $c\cdot y \le \rho$. Discharging $\leaf$ means proving this system infeasible.

\paragraph{Untrusted \scip{}, checked \vipr{}.} We hand the system to \scip-Ex, which
returns a \vipr{} certificate of infeasibility. Our checker replays the \code{DER}
list over $\Q$: for each \code{lin} step with multipliers $y$ it verifies
sign-consistency and that the aggregate dominates the stated constraint; \code{rnd}
adds an integer floor/ceil under integrality; \code{uns} combines two children over a
complementary integer split and discharges the branch assumption. Acceptance requires
the final derived constraint to be an absurdity with empty assumption set. The only
non-arithmetic fact used is the soundness (\emph{easy}) direction of Farkas:

\begin{lemma}[Farkas, checking direction]\label{lem:farkas}
For $A\in\Q^{m\times n}$, $b\in\Q^m$, $y\in\Q^m$: if $y\ge 0$, $y^{\!\top}\!A = 0$ and
$y^{\!\top} b < 0$, then $\{x : Ax\le b\}=\varnothing$.
\end{lemma}
\begin{proof}
If $x$ were feasible, $0=(y^{\!\top}\!A)x=y^{\!\top}(Ax)\le y^{\!\top} b<0$, using
$y\ge 0$ and $Ax\le b$ termwise. \qed
\end{proof}

Lemma~\ref{lem:farkas} is a two-line argument from $\Q$'s ordered-field structure and
finite sums, entirely within \mathlib. Notably we do \emph{not} need the hard
(existence) direction of Farkas---\mathlib's general hyperplane-separation
result~\cite{mathlib}---for soundness; it is required only to prove \tool{}
\emph{complete} (that a true query always admits a certificate), which we leave to
future work.

\paragraph{Correspondence.} A \vipr{} file is self-contained, so we additionally
check that its \code{VAR}/\code{CON}/\code{OBJ} sections \emph{coincide} with the system
$\tool$ built from the \code{.net} file, the box, and the leaf's splits---otherwise
\scip{} could certify a different problem. This byte-level correspondence is checked
when assembling the command-line pipeline (\S\ref{sec:implementation}) and is what
links Lemma~\ref{lem:farkas}'s conclusion back to $\reg(\leaf)$.

\section{Formalization in \lean{}~4}\label{sec:formalization}

\paragraph{Exact substrate.} Everything is over $\Q$ (\mathlib's \code{Rat}), a
computable linearly ordered field with decidable comparison. Network weights arrive as
IEEE-754 bit patterns; each finite float is exactly the dyadic rational
$(-1)^s\, m\, 2^{e-\text{bias}-p}$, and we decode bits to $\Q$ with \emph{zero}
rounding (a $\sim$30-line function), so the \lean{} network is bit-identical to the
one the verifier reasoned about. No step of the checker uses \code{Float}.

\paragraph{Network semantics.} \code{Model/Network} defines the layer types and the
denotational forward map $\net:\ (\code{Fin }n_0\to\Q)\to(\code{Fin }n_\ell\to\Q)$;
$\relu$ is $\max(\cdot,0)$, exact over $\Q$.

\paragraph{Encoding soundness.} \code{Model/Encoding} builds the rational system and
we show it \emph{over-approximates} the real behaviour (feasible set $\supseteq$ the
graph of $\net$ on the region). Two cornerstones are mechanized and axiom-clean:
\begin{lemma}[Big-$M$ ReLU, \code{reluBigM\_sound}]\label{lem:relu}
For $lb\le t\le ub$, the exact value $v=\max(t,0)$ and binary $a=[\,t\ge 0\,]$ satisfy
$v\le t-lb(1-a)$, $t\le v$, $v\le ub\,a$, and $0\le v$.
\end{lemma}
\begin{lemma}[Interval bounds, \code{affine\_interval\_sound}]\label{lem:ivl}
For every $x$ in the box, $\sum_j(w_j^{+}lo_j+w_j^{-}hi_j)+b\le \sum_j w_j x_j+b\le
\sum_j(w_j^{+}hi_j+w_j^{-}lo_j)+b$, with $w^{+}=\max(w,0)$, $w^{-}=\min(w,0)$.
\end{lemma}
Lemma~\ref{lem:relu} says each ReLU gadget admits the real activation as a feasible
point; Lemma~\ref{lem:ivl} validates the big-$M$ constants and the stability
decisions (\code{stable\_active}/\code{stable\_inactive}). Threading these through the
layers---choosing each binary from the true neuron sign---gives the assembly, the
central theorem (mechanized for arbitrary-depth MLPs; Conv/CNN is future work; see
\S\ref{sec:formalization} ``Mechanization status''):
\begin{lemma}[Encoding over-approximation]\label{lem:enc}
For every leaf $\leaf$ and every $x\in\reg(\leaf)$ there is an assignment $z$ to the
encoded variables with $A_\leaf z \le b_\leaf$ and terminal value $c\cdot\net(x)$.
\end{lemma}
Combining Lemmas~\ref{lem:farkas} and~\ref{lem:enc}: if the certificate proves
$A_\leaf x\le b_\leaf \wedge c\cdot y\le\rho$ infeasible, then no $x\in\reg(\leaf)$ has
$c\cdot\net(x)\le\rho$, which is Lemma~\ref{lem:leaf}.

\begin{figure}[t]
\begin{lstlisting}
theorem certified_sound
    (net : Network) (box : Box net.inDim) (obj : Objective net.outDim)
    (aptp : Aptp) (certs : LeafKey -> Vipr) :
    check net box obj aptp certs = .certified ->
    forall x, x in box -> obj.c dot net.eval x > obj.rhs := by
  intro hchk x hx
  obtain <L, hLmem, hxL> := coverage_sound  (aptp_coverage_ok hchk) x hx
  exact leaf_refuted (leaf_cert_ok hchk L hLmem) x hxL
\end{lstlisting}
\caption{The target top-level theorem (elided), composing the mechanized coverage
check (Lemma~\ref{lem:cover}) and per-leaf refutation (Lemma~\ref{lem:leaf}).}
\label{fig:thm}
\end{figure}

\paragraph{Reflective checkers.} Each check---Farkas step, rounding, unsplit,
coverage---is a \code{Bool}-valued function proved to imply its semantic
consequence (\emph{soundness lemma}), then \emph{evaluated by the kernel} on concrete
rationals. We deliberately avoid \code{native\_decide}: acceptance is a kernel
reduction, so the compiler and runtime stay out of the trusted base.

\paragraph{Mechanization status.} Everything reported in this section is
machine-checked and \emph{axiom-clean} (depending only on \code{propext},
\code{Classical.choice}, \code{Quot.sound}---no \code{sorry}, no \code{native\_decide}).
Foundationally: the Farkas checking direction (Lemma~\ref{lem:farkas},
\code{farkas\_infeasible}), the coverage check (Lemma~\ref{lem:cover},
\code{coverage\_sound}, with a complete Boolean-cube enumeration proof), and the two
encoding cornerstones (Lemmas~\ref{lem:relu} and~\ref{lem:ivl}); the fuller results
(arbitrary-depth encoding, the certificate checker, and both parsers) follow. The
exact substrate is validated operationally: the $\Q$
forward evaluator agrees \emph{bit-for-bit} with an independent rational forward pass
on the running network, and the `.aptp`/`.net` parsers reproduce the reference
Python parser's output on the sample proofs.

\emph{Encoding soundness} (Lemma~\ref{lem:enc}) is proved for the full
\textbf{arbitrary-depth MLP} class $\mathrm{Lin}\!\to\!\relu\!\to\!\cdots\!\to\!\mathrm{Lin}$
(a \code{Linear}, then any number of $(\relu,\mathrm{Linear})$ blocks), for the
\textbf{leaf-aware} encoder \code{MLP.encFold} the tool actually runs: a neuron the
leaf sign-fixes (or the intervals stabilize) is folded to $\mathrm{post}=\mathrm{pre}$
or $\mathrm{post}=0$ with \emph{no} integer variable, and only genuinely-unstable
neurons get the four big-$M$ rows and a binary. By induction on depth,
\code{encFold\_overapprox} builds the real trace and shows it satisfies every emitted
row (box, every affine layer, the per-neuron fold or big-$M$ rows, and the leaf's sign
rows) with the objective evaluating to $c\cdot\net(x)$; \code{encFold\_binIds\_int}
shows the trace is $\{0,1\}$-valued on the (only) declared integer variables;
\code{certified\_sound\_mlp\_fold} composes with \code{refute\_of\_cert} for per-leaf
refutation, and the \code{satLeaf}-phrased variant slots directly into
\code{certified\_sound\_abstract} (Fig.~\ref{fig:thm}). There is exactly \emph{one}
encoder, and it is the one both run and proven. All axiom-clean. (The soundness
statement is on a dimension-indexed \code{MLP} normal form ending in a linear scoring
layer; \code{Flatten} is the identity and elided, consecutive linears/ReLUs collapse,
and \emph{Conv/CNN is out of scope}, left as future work.)

\emph{Certificate replay} is mechanized as a complete, automatic, runnable checker
\code{checkSem} over the certificate's own \emph{sensed} constraints
($\le$/$\ge$/$=$), implementing full VIPR~v1.0 verification semantics and consuming
the solver's flat, index-referenced derivation list \emph{directly} --- no
tree-reshaping step exists. Its ingredients are proved axiom-clean: the
\emph{suitable linear combination} (signed multipliers over mixed senses, per-term
monotonicity generalizing Lemma~\ref{lem:farkas}), constraint \emph{domination}
including the spec's absurdity-dominates-anything rule, \code{rnd} (integer rounding
of the combined row, both directions), and \code{uns} (an integer case-split: the two
discharged rows must be complementary integer bounds $f\le\beta\,/\,f\ge\beta{+}1$,
each branch's row must dominate the stated row, and each derived row carries the set
of still-open assumptions, discharged exactly at its split). The replay invariant
(\code{sfreplay\_holds}: every derived row holds under its open assumptions) gives
\code{checkSemCore\_sound}: if the checker accepts --- the final row an absurdity with
no open assumptions --- no valuation with the declared variables integral satisfies
the \code{CON} rows, i.e.\ the MILP is infeasible. \code{checkSem\_sound} composes
this with the VIPR parser. The checker is validated on \emph{real} certificates from
official exact \scip{}~10.0.1: signed-multiplier and Gomory-cut proofs on the worked
example's leaves, and a genuinely branched proof (\code{asm}/\code{uns}) on a knapsack
instance --- each also correctly \emph{rejected} when the integrality declarations are
dropped.

\emph{Parser verification.} The `.net` reader is \textbf{fully verified}:
\code{parseNet\_printNet} proves \code{parseNet (printNet r) = .ok (rawToNet r)} for
well-formed `r` (axiom-clean), composing stage round-trips for the float32-bit weight
decode (\code{parseWeight\_printHex}), the weight list (\code{takeWeights\_printHex}),
the decimal dimensions, the tokenizer (\code{tokenize\_jn}), and the statement reader
(\code{readStatements\_single})---each proved to invert a corresponding printer. This
required making the layer loop \emph{total}: the original \code{partial def} is opaque
to the kernel (no equational lemmas), so we rewrote it as a fuel-structural recursion
(\code{parseLayersFuel}, fuel $=$ token count $+1$), behaviour-preserving
(regression-checked: \code{parseNet} yields byte-identical exact outputs on the
sample). The \code{.aptp} proof-tree reader is verified the same way
(\code{parseAptp\_printAptp}), which additionally required de-\code{partial}-izing the
S-expression parser and rewriting the two-pass \code{parseAptp} loops as total
\code{List.foldl}s (behaviour-preserving; regression-checked on both sample proofs).
\textbf{Both} parsers are thus out of the trusted base.

The soundness is also bridged to the \emph{executable} side: \code{toMLP} converts a
parsed \code{Array}-based \code{Network} (the MLP normal form) to the
dimension-indexed \code{MLP}, \code{toMLP\_eval} proves \code{Network.eval} agrees with
\code{MLP.eval}, and \code{certified\_sound\_network\_fold} transports the per-leaf
refutation onto the parsed network's own output---so the runnable network is exactly
the one the theorem is about (all axiom-clean).

\emph{End-to-end capstone.} The pieces compose into one theorem,
\code{certify\_network\_sound}: if the coverage check passes and every leaf's
\code{leafCheck} passes --- \code{checkSem} accepts the certificate, every certificate
\code{CON} row semantically matches a row of that leaf's \code{encFold} encoding (or a
binary's $[0,1]$ bound), and the declared integer variables are the encoder's binaries
--- then $\forall x\in\dom,\ c\cdot\net(x)>\rho$ on the parsed network. The
correspondence is checked up to the solver's variable renumbering (undone, untrusted,
from the certificate's own variable names) and term reordering (the proven
order-insensitive form equality); the true-sign vector \code{MLP.trueSign} with
\code{signMatch\_trueSign} supplies the coverage side. The CLI computes exactly this
hypothesis, so its \sound{} verdict carries the theorem's conclusion; trusted beyond
Lean's kernel are only the OS/file-I/O glue and \code{main}'s conjunction of the
per-leaf verdicts.

Remaining, and none requiring trust in \scip{} or floats: Conv/CNN encoding (future
work).

\paragraph{Trusted computing base.} The TCB is: the \lean{} kernel and $\Q$
arithmetic, and the \code{.net} file faithfully representing the intended network (its
generation is out of scope---the tool reads only the exact format). \textbf{Both}
input parsers---the network file and the proof-tree file---are now \emph{fully
verified} (\code{parseNet\_printNet}, \code{parseAptp\_printAptp}) and out of the TCB.
\scip, SoPlex, the \vipr{} certificate, and any certificate-completion tooling are
\emph{not} trusted.

\section{Design rationale: why each step is needed}\label{sec:rationale}

Each ingredient of \S\ref{sec:checking}--\ref{sec:formalization} is forced by a
specific soundness or tractability concern; we spell these out because they explain
both the architecture and where the proof effort concentrates.

\begin{description}
  \item[Why exact rationals, not floats.] The artifact being built is a
    \emph{soundness} checker; if it computed in floating point, a rounding error
    could make it accept a false proof---exactly the failure mode we are trying to
    rule out in the underlying verifier. Since every finite binary32 weight
    \emph{is} a dyadic rational, decoding bits to $\Q$ (\S\ref{sec:formalization})
    loses nothing, so exactness costs only speed, never faithfulness.

  \item[Why SCIP/\vipr{}, not Gurobi.] A solver that returns only ``infeasible''
    forces us to \emph{trust} it, re-introducing a large floating-point dependency
    in the trusted base. An exact solver that emits a \vipr{} certificate lets us
    move it \emph{out} of the trusted base: we keep the answer only if we can
    re-derive it.

  \item[Why re-check, not re-solve.] Solving an MILP inside the kernel would require
    a verified branch-and-cut engine---enormous, and slow in kernel reduction.
    Checking exploits the search/verification asymmetry: a \vipr{} \code{lin} step is
    validated by the \emph{easy} direction of Farkas (\code{farkas\_infeasible},
    \code{farkas\_le}), a two-line contradiction (\S\ref{sec:certificates}). We
    therefore invest in a certificate replayer, not a solver.

  \item[Why two obligations (refutation \emph{and} coverage).] The goal
    \eqref{eq:property} is universally quantified over the box. A proof tiles the box
    into activation regions; refuting each region is vacuous unless the regions
    \emph{cover} the box, and coverage is vacuous unless each region is refuted.
    Neither obligation alone is sound, which is why
    Theorem~\ref{thm:main} needs both Lemmas~\ref{lem:leaf} and~\ref{lem:cover}.

  \item[Why coverage is a Boolean-cube tautology (an over-approximation).] The
    \emph{reachable} sign patterns are as hard to characterize as the verification
    problem itself. Checking that the leaves cover the entire free cube
    $\{0,1\}^{\#\text{neurons}}$ is decidable and \emph{sufficient} (every reachable
    pattern is some cube point), at the cost of rejecting proofs that cover only
    reachable patterns---a completeness, not soundness, concession
    (\S\ref{sec:checking}). Getting the polarity-flip right
    (\code{coverage\_sound} negates literals) is what makes the enumeration
    \emph{provably} equivalent to non-tautology.

  \item[Why leaves are MILPs, not LPs.] Partial patterns leave unstable neurons
    unbranched, so their binary indicators survive into each leaf. A per-leaf
    certificate is thus a full \vipr{} tree (\code{lin}+\code{rnd}+\code{uns}), and
    the checker's constraint substrate must be the sparse, variable-indexed
    \code{Le}/valuation form (\code{Cert/LinCon}) that certificates actually speak,
    not a single dense multiplier vector. The affine \emph{complete-pattern} case
    (\code{certified\_sound\_singleLinear}) collapses to one \code{lin} step and is
    the natural first milestone.

  \item[Why the over-approximation \emph{direction} matters.] Soundness needs the
    encoded feasible set to \emph{contain} the network's real behaviour: then
    infeasibility of ``encoding $\wedge$ $c\cdot y\le\rho$'' rules out a real
    counterexample. This is why \code{reluBigM\_sound} shows the real activation
    \emph{satisfies} the big-$M$ gadget (feasible), rather than proving the gadget
    \emph{equals} the ReLU. A looser encoding is harmless; a tighter-than-reality one
    would be unsound---hence also why the interval bounds must be \emph{validated}
    (\code{affine\_interval\_sound}) and any folded ``stable'' neuron re-verified
    (\code{stable\_active}/\code{stable\_inactive}) rather than trusted from the
    producer.

  \item[Why cornerstones first, then assembly.] The encoding-soundness theorem is a
    layer induction with variable bookkeeping. Isolating the per-neuron facts as
    reusable, independently kernel-checked lemmas (Lemmas~\ref{lem:relu}
    and~\ref{lem:ivl}) turns that induction into pure plumbing: each layer is
    discharged by a cornerstone, and the objective row by \code{fuse\_linear}.

  \item[Why an abstract composition interface.] \code{certified\_sound\_abstract}
    and \code{refute\_of\_cert} prove the top-level structure
    (coverage $\oplus$ per-leaf refutation $\Rightarrow$ property) against the
    \emph{real} \code{coverage\_sound} and Farkas core, parameterized by the two
    pieces still under construction (the sign map from the trace, and the encoded
    rows). This lets the finished parts be kernel-checked now and the encoding
    assembly slot in without touching the logical glue.

  \item[Why kernel checking, not \code{native\_decide}.] \code{native\_decide} would
    add the Lean compiler and runtime to the trusted base---unacceptable for a tool
    whose entire purpose is a small trusted base. Every check is a reflective
    \code{Bool} function evaluated by the kernel, so acceptance is a kernel reduction.
    (The one representation friction---the executable network uses \code{Array} while
    proofs use \code{Fin}/\code{Finset} sums---is bridged with \code{List.ofFn} and
    \code{sum\_ofFn}, keeping the assembly proof kernel-friendly.)
\end{description}

\section{Implementation}\label{sec:implementation}

\tool{} is a \lean{}~4 package ($\approx$5.3k lines) built on \mathlib. Implemented and
kernel-checked: the exact numeric layer (\code{Numeric/Float}); the S-expression lexer
and the \code{.aptp}/\code{.net}/\vipr{} parsers (\code{Ast/*}) with full round-trip
proofs (\code{Ast/NetRoundtrip}, \code{Ast/AptpRoundtrip}); the exact network evaluator
(\code{Model/Network}); the encoding cornerstones and the executable encoder
(\code{Model/Encoding}, \code{Model/Encoder}); the \emph{arbitrary-depth} MLP
encoding-soundness proof (\code{Model/EncodingSound}); the Farkas checks
(\code{Cert/LinComb}, \code{Cert/LinCon}); the coverage check
(\code{Coverage/Tautology}); the complete proof-by-cases certificate checker
(\code{Cert/Vipr}: \code{lin}/\code{rnd}/\code{uns} soundness, \code{RefTree}, the
runnable \code{checkRefTree}); the \code{Network}$\to$\code{MLP} bridge transporting soundness onto the executable
network (\code{Model/NetworkMLP}); and the composition scaffolding plus the VIPR
endpoint (\code{Pipeline/Compose}, \code{Pipeline/Affine}, \code{Pipeline/ViprCheck}).
The command-line driver \code{aptpcheck N.net P.aptp [-{}-vipr-dir D/]} is assembled and
runs end to end: it invokes exact \scip{} per leaf and validates each returned
certificate with the verified \code{leafCheck}, so its \sound{} verdict is the
hypothesis of \code{certify\_network\_sound} (\code{Pipeline/EndToEnd}); there is no
flat-to-tree conversion --- the checker consumes the solver's derivation list directly.
Under construction: Conv/CNN. The proof format and neuron numbering match \pytool, so
its proofs are consumed unchanged.

\section{Evaluation}\label{sec:eval}

\paragraph{End-to-end on the worked example (real, reproducible).} We ran the
assembled \code{aptpcheck} on \code{sample.net}/\code{sample.aptp} against official
\code{scipoptsuite-10.0.1} (exact rational mode, exact \textsc{SoPlex}). The tool emits
each leaf as an exact-rational \code{.mps} from the verified leaf-aware encoder, exact
\scip{} returns a \vipr{} certificate, and the kernel-checked \code{leafCheck} validates
the certificate \emph{and} its correspondence to the leaf's encoding. All four leaves
are accepted and the tool reports \sound{} (Step~10, \S\ref{sec:example}).
Table~\ref{tab:eval} gives the real per-leaf data. The integer-variable column equals
the number of neurons the leaf leaves unfixed --- the leaf-aware encoding at work. Three
certificate families appear: signed combinations (\code{lin}), Gomory rounding cuts
(\code{rnd}), and --- in the branching stress test --- assumptions and a case-split
(\code{asm}/\code{uns}). Every certificate is genuine, untrusted solver output, re-checked
in exact~$\Q$ by the verified checker.

\begin{table}[t]
\centering
\caption{Real end-to-end run of \tool{} on the worked example with official exact
\scip{} 10.0.1 (per-leaf MILPs from the verified leaf-aware encoder; each certificate
re-checked by the kernel-checked \code{leafCheck}). Bottom: a branching stress test
(exact \scip{} solves it by branch-and-bound; the certificate is accepted by \code{checkSem}).}
\label{tab:eval}
\begin{tabular}{lrrrrl}
\toprule
Instance & \#vars & \#int & \#constraints & \#\vipr{} steps & reasons used \\
\midrule
$\leaf_1=[-4]$      & 16 & 3 & 38 & 97  & \code{lin} \\
$\leaf_2=[-2,4]$    & 16 & 2 & 35 & 273 & \code{lin}, \code{rnd}$\times2$ \\
$\leaf_3=[2,1,4]$   & 16 & 1 & 32 & 177 & \code{lin} \\
$\leaf_4=[2,-1,4]$  & 16 & 1 & 32 & 983 & \code{lin}, \code{rnd}$\times1$ \\
\midrule
knapsack $6x{+}10y{+}15z{=}23$ & 3 & 3 & 8 & 32 &
  \code{asm}$\times2$, \code{lin}$\times21$, \code{rnd}$\times8$, \code{uns}$\times1$ \\
\midrule
\multicolumn{6}{l}{\emph{Verdict (sample):} coverage OK; all leaves \sound{}
$\Rightarrow \forall x\in\dom.\ Y_0>Y_1$ (\code{certify\_network\_sound}).} \\
\bottomrule
\end{tabular}
\end{table}

\paragraph{Scaling (future work).} A broader evaluation on VNN-COMP benchmarks
(ACAS-Xu, MNIST-FC) will quantify (Q1) agreement with the reference floating-point
checker, (Q2) the overhead of exact arithmetic and kernel-level replay, and (Q3) how
\vipr{} certificate size and rational-coefficient bit-width grow with network size. The
worked example already stresses the full chain (exact-rational MPS, real signed-multiplier
and Gomory-cut certificates, kernel-checked replay); scaling it is an engineering effort
we leave to future work.

\section{Related work}\label{sec:related}

\paragraph{NNV and proof production.} Complete verifiers such as
Reluplex/Marabou~\cite{katz2017reluplex} and the CROWN/$\beta$-CROWN
family~\cite{wang2021betacrown} dominate VNN-COMP~\cite{vnncomp2023}. Marabou can
\emph{produce proofs}: a proof tree with Farkas coefficients at LP leaves and lemmas
for \relu{} splits~\cite{isac2022marabou}, checked by a dedicated (unverified)
checker. The APTP format generalizes this to a verifier-independent artifact. \tool{}
differs by reducing checking to a \emph{proof-assistant kernel} and by delegating
arithmetic to an untrusted \emph{certificate-producing} solver.

\paragraph{Verified checking of solver results.} The ``untrusted oracle + reflective
verified checker'' pattern is standard in Coq \code{micromega}/\code{lra} for linear
arithmetic~\cite{besson2007micromega} and underlies verified SAT/LP checkers. For
integer programming, \vipr~\cite{cheung2017vipr} and exact
\scip~\cite{eifler2022exact} provide certificates checked by \code{viprchk}, which is
trusted by inspection; we re-implement and \emph{verify} such a checker in \lean.

\paragraph{Verified neural-network reasoning.} Prior efforts formalize robustness or
network semantics in Coq/Isabelle and check certificates in
Imandra~\cite{desmartin2022checkinn}. We are, to our knowledge, the first to give a
kernel-checked, exact soundness formalization---encoding, coverage, certificate
checking, and both parsers all machine-checked---for a \emph{verifier-independent}
NNV proof format backed by an untrusted MILP solver.

\section{Conclusion}\label{sec:conclusion}

We presented \tool, an APTP proof checker formalized in \lean/\mathlib{} over exact
rationals, trusting neither floating point nor any MILP solver: \scip{} runs untrusted
and its \vipr{} certificate is re-checked in the kernel. The guarantee factors cleanly
into per-leaf refutation, coverage, and encoding soundness. Machine-checked and
axiom-clean: encoding soundness for arbitrary-depth ReLU MLPs, a complete automatic
proof-by-cases \vipr{} certificate checker (Farkas/rounding/split), the coverage
check, both input parsers by round-trip (out of the trusted base), and a bridge that
transports the soundness onto the executable, parsed network. What remains is
engineering, not new trust: assembling the command-line driver (with an untrusted
flat-\vipr{}-to-tree converter), and extending to convolutional and residual
architectures at scale. We then target \emph{completeness} (via \mathlib's Farkas/LP
duality) and parallel certificate checking.

\bibliographystyle{splncs04}
\bibliography{references}

\end{document}
